\begin{lstlisting}[caption={Go script for performance monitoring via perf for Prometheus}, label={lst:perf-prometheus-exporter}]
package main
import (
	"fmt"
	"io/ioutil"
	"net/http"
	"os/exec"
	"strconv"
	"strings"
	"time"

	"gopkg.in/yaml.v2"
)

// Config represents the structure of the YAML configuration file.
// It includes high-level performance counters (HCP), software events (SW), and hardware events (HW).
type Config struct {
	HCP []string `yaml:"hcp"`
	SW  []string `yaml:"sw"`
	HW  []string `yaml:"hw"`
}

// Global variable to store the most recent performance data.
var perfData string

// getPerfData runs `perf stat` with the specified events for a short duration (0.1s)
// and returns the raw output as a string.
func getPerfData(events []string) (string, error) {
	eventList := strings.Join(events, ",")
	cmd := exec.Command("perf", "stat", "-e", eventList, "--", "sleep", "0.1")
	output, err := cmd.CombinedOutput()
	if err != nil {
		return "", err
	}
	return string(output), nil
}

// parsePerfData filters and extracts only the relevant lines from perf output,
// based on the provided list of events.
func parsePerfData(data string, events []string) string {
	lines := strings.Split(data, "\n")
	var result []string
	for _, line := range lines {
		for _, event := range events {
			if strings.Contains(line, event) {
				result = append(result, line)
				break
			}
		}
	}
	return strings.Join(result, "\n")
}

// formatForPrometheus converts the filtered performance data into Prometheus-compatible metrics format.
func formatForPrometheus(data string) string {
	lines := strings.Split(data, "\n")
	var result []string
	for _, line := range lines {
		parts := strings.Fields(line)
		if len(parts) >= 2 {
			// Clean up the metric name
			metricName := strings.ReplaceAll(parts[1], ".", "_")
			// metricName = strings.ReplaceAll(metricName, "-", "_")
			// Remove decimal points from the number and parse it
			metricValue := strings.ReplaceAll(parts[0], ".", "")
			if value, err := strconv.ParseFloat(metricValue, 64); err == nil {
				// Format value appropriately for Prometheus (integer or float)
				if value == float64(int(value)) {
					result = append(result, fmt.Sprintf("%s %d", metricName, int(value)))
				} else {
					result = append(result, fmt.Sprintf("%s %f", metricName, value))
				}
			}
		}
	}
	return strings.Join(result, "\n")
}

// metricsHandler handles HTTP requests and returns the formatted performance metrics.
func metricsHandler(w http.ResponseWriter, r *http.Request) {
	fmt.Fprintf(w, formatForPrometheus(perfData))
}

// chunkEvents splits the list of events into smaller chunks of a given size.
// This helps avoid issues with perf's limit on the number of simultaneous counters.
func chunkEvents(events []string, chunkSize int) [][]string {
	var chunks [][]string
	for i := 0; i < len(events); i += chunkSize {
		end := i + chunkSize
		if end > len(events) {
			end = len(events)
		}
		chunks = append(chunks, events[i:end])
	}
	return chunks
}

// validateEvent checks if a given performance event is available on the system by using `perf list`.
func validateEvent(event string) bool {
	cmd := exec.Command("perf", "list")
	output, err := cmd.CombinedOutput()
	if err != nil {
		fmt.Println("Error listing perf events:", err)
		return false
	}
	return strings.Contains(string(output), event)
}

func main() {
	// Load the YAML configuration file
	config := Config{}
	data, err := ioutil.ReadFile("config.yml")
	if err != nil {
		fmt.Println("Error reading config file:", err)
		return
	}
	err = yaml.Unmarshal(data, &config)
	if err != nil {
		fmt.Println("Error parsing config file:", err)
		return
	}

	// Validate and collect all available events from the config
	var validEvents []string
	for _, event := range append(config.HCP, config.SW...) {
		if validateEvent(event) {
			validEvents = append(validEvents, event)
		}
	}
	for _, event := range config.HW {
		if validateEvent(event) {
			validEvents = append(validEvents, event)
		}
	}

	// Split the events into manageable chunks (TODO: dynamically adjust based on CPU performance)
	eventChunks := chunkEvents(validEvents, 6)

	// Start a background goroutine to periodically collect performance data
	go func() {
		for {
			var allPerfData []string
			for _, chunk := range eventChunks {
				data, err := getPerfData(chunk)
				if err != nil {
					fmt.Println("Error:", err)
					allPerfData = append(allPerfData, "Error fetching data")
				} else {
					allPerfData = append(allPerfData, parsePerfData(data, chunk))
				}
			}
			// Combine data from all chunks into one string for Prometheus export
			perfData = strings.Join(allPerfData, "\n")
			time.Sleep(100 * time.Millisecond)
		}
	}()

	// Start the HTTP server that exposes the `/metrics` endpoint for Prometheus scraping
	http.HandleFunc("/metrics", metricsHandler)
	fmt.Println("Starting server on :9011")
	if err := http.ListenAndServe(":9011", nil); err != nil {
		fmt.Println("Error starting server:", err)
	}
}
\end{lstlisting}
